\section*{Conclusion}

We presented \ours{}, a benchmark of 63{,}000 clinical QA pairs over MIMIC-IV covering 63 concepts stratified by compositional reasoning depth, designed to evaluate whether LLMs can generate modular, composable libraries for EHR data processing. As a strong baseline, we proposed clinical autoformalization, which converts natural-language clinical guidelines into verified Python libraries via a QA-driven refinement loop, removing the need for manual tool engineering. The resulting libraries generalize across patient populations, EHR databases, and evolving clinical guidelines while reducing per-query token usage compared to zero-shot approaches. We see opportunity in extending \ours{} to multimodal data sources and in using the verified libraries as supervised signal for training smaller, deployable clinical agents.

\section*{Limitations}

Several limitations bound the current work. First, the verification set used during autoformalization is small (10\% of each concept's QA pairs, approximately 100 examples), which may allow rare edge cases---particularly unusual lab values or atypical patient trajectories---to escape detection. Second, the benchmark is grounded in MIMIC-IV, a dataset from a single US academic medical center, and the autoformalized functions inherit the clinical conventions of that institution; generalization to EHR systems with different schema designs, coding practices, or patient populations is not guaranteed without re-running the autoformalization loop. Third, the natural-language guidelines used as input are generated by an LLM from PubMed abstracts, introducing the possibility of clinical inaccuracies that propagate into the verified functions. Fourth, the longitudinal reasoning component of the benchmark is still under development and is not yet evaluated in the current experiments.

\section*{Broader Impacts}

This work has potential benefits for clinical informatics by reducing the expert effort required to deploy LLM-based EHR analysis tools and by promoting reproducible, verifiable code over ad-hoc zero-shot generation. A library of QA-verified clinical functions is more auditable than a sequence of one-off model outputs, which is relevant for safety-critical applications. At the same time, automated extraction of clinical concepts from EHR data carries risks if deployed without appropriate clinical validation: errors in autoformalized functions could propagate silently to clinical decision support systems. Additionally, MIMIC-IV is demographically skewed toward patients treated at Beth Israel Deaconess Medical Center in Boston, and function libraries trained on this distribution may perform less reliably for underrepresented patient groups. We recommend that any downstream deployment of autoformalized libraries be subject to institution-specific validation before clinical use.
